
\begin{tikzpicture}[
    scale=0.75, transform shape,
    % Node styles
    input_node/.style={circle, draw=black!80, fill=green!30, minimum size=8mm},
    conv_block/.style={rectangle, draw=black!80, fill=blue!20, minimum width=25mm, minimum height=15mm, rounded corners=2mm},
    temporal_block/.style={rectangle, draw=black!80, fill=orange!20, minimum width=40mm, minimum height=50mm, rounded corners=3mm},
    fc_layer/.style={rectangle, draw=black!80, fill=purple!20, minimum width=30mm, minimum height=12mm, rounded corners=2mm},
    output_node/.style={circle, draw=black!80, fill=red!30, minimum size=8mm},
    % Arrow styles
    data_flow/.style={-latex, thick, black!70},
    skip_connection/.style={-latex, thick, blue!60, dashed},
    % Labels
    label_style/.style={font=\footnotesize\sffamily},
    block_label/.style={font=\small\bfseries\sffamily}
]

% Input Layer
\node[block_label] at (-2, 0) {Input};
\foreach \i in {1,...,5} {
    \node[input_node] (input\i) at (0, -\i*0.8 + 2.4) {\tiny\i};
}

% Temporal Block 1 (dilation=1)
\begin{scope}[shift={(3,0)}]
    \node[temporal_block, minimum height=55mm] (tb1) at (0,0) {};
    \node[block_label] at (0, 2.5) {Temporal Block 1};
    \node[label_style] at (0, 2) {dilation = 1};
    
    % Internal structure
    \node[conv_block] (conv1_1) at (0, 1.2) {Conv1d(5→64)};
    \node[label_style, below=2pt of conv1_1] {k=2, p=1};
    
    \node[conv_block] (conv1_2) at (0, -0.5) {Conv1d(64→64)};
    \node[label_style, below=2pt of conv1_2] {k=2, p=1};
    
    % Skip connection
    \node[conv_block, minimum width=20mm, minimum height=8mm] (skip1) at (0, -2) {1×1 Conv(5→64)};
    
    % Operations
    \node[label_style, right=2pt of conv1_1] {+ReLU+Drop};
    \node[label_style, right=2pt of conv1_2] {+ReLU+Drop};
\end{scope}

% Temporal Block 2 (dilation=2)
\begin{scope}[shift={(8,0)}]
    \node[temporal_block, minimum height=55mm] (tb2) at (0,0) {};
    \node[block_label] at (0, 2.5) {Temporal Block 2};
    \node[label_style] at (0, 2) {dilation = 2};
    
    % Internal structure
    \node[conv_block] (conv2_1) at (0, 1.2) {Conv1d(64→64)};
    \node[label_style, below=2pt of conv2_1] {k=2, p=2, d=2};
    
    \node[conv_block] (conv2_2) at (0, -0.5) {Conv1d(64→64)};
    \node[label_style, below=2pt of conv2_2] {k=2, p=2, d=2};
    
    % Operations
    \node[label_style, right=2pt of conv2_1] {+ReLU+Drop};
    \node[label_style, right=2pt of conv2_2] {+ReLU+Drop};
\end{scope}

% Temporal Block 3 (dilation=4)
\begin{scope}[shift={(13,0)}]
    \node[temporal_block, minimum height=55mm] (tb3) at (0,0) {};
    \node[block_label] at (0, 2.5) {Temporal Block 3};
    \node[label_style] at (0, 2) {dilation = 4};
    
    % Internal structure
    \node[conv_block] (conv3_1) at (0, 1.2) {Conv1d(64→64)};
    \node[label_style, below=2pt of conv3_1] {k=2, p=4, d=4};
    
    \node[conv_block] (conv3_2) at (0, -0.5) {Conv1d(64→64)};
    \node[label_style, below=2pt of conv3_2] {k=2, p=4, d=4};
    
    % Operations
    \node[label_style, right=2pt of conv3_1] {+ReLU+Drop};
    \node[label_style, right=2pt of conv3_2] {+ReLU+Drop};
\end{scope}

% FC Layer
\node[fc_layer] (fc) at (18, 0) {Linear(64→5)};
\node[block_label] at (18, 1) {Output Layer};

% Output nodes
\node[block_label] at (21, 0) {Output};
\foreach \i in {1,...,5} {
    \node[output_node] (output\i) at (20.5, -\i*0.8 + 2.4) {\tiny\i};
}

% Connections
% Input to TB1
\foreach \i in {1,...,5} {
    \draw[data_flow] (input\i) -- (tb1.west);
}

% TB1 to TB2
\draw[data_flow] (tb1) -- (tb2);

% TB2 to TB3
\draw[data_flow] (tb2) -- (tb3);

% TB3 to FC
\draw[data_flow] (tb3) -- (fc);

% FC to Output
\foreach \i in {1,...,5} {
    \draw[data_flow] (fc.east) -- (output\i);
}

% Add receptive field visualization
\begin{scope}[on background layer]
    \node[draw=gray!30, fill=gray!10, rounded corners=5pt, 
          fit=(input1)(input5), inner sep=3pt, label=below:{\footnotesize t-7...t}] {};
    
    % Receptive field indicators
    % \draw[<->, thick, gray!60] ([yshift=-4cm]input1.south) -- 
    %     node[midway, below, font=\footnotesize] {Receptive Field = 8} 
    %     ([yshift=-4cm]input5.south -| tb3.east);
\end{scope}

% Add detailed annotations
\node[align=center, font=\footnotesize] at (10, -4) {
    Causal padding ensures\\output at time $t$ only\\depends on $t$ and earlier
};

% Title
\node[font=\Large\bfseries] at (9, 4) {Temporal Convolutional Network (TCN)};
\node[font=\normalsize] at (9, 3.3) {3 Blocks, Exponentially Increasing Dilation};

\end{tikzpicture}
